% --- Archon physics formalization source begin ---
% archon:physics
% archon:covers PhyXMiniProblems/problem_phyx_mini_0328.lean
% archon:source-report reports/phyx_mini/problem_phyx_mini_0328.source.json

\chapter{Physics problem phyx\_mini\_0328}
\label{ch:PhyXMiniProblems_problem_phyx_mini_0328}

\paragraph{Problem source.}
\#\# Physical scenario

In figure, a rod \$R\$ is clamped at its center; a disk \$D\$ at its end projects into a glass tube that has cork filings spread over its interior. A plunger \$P\$ is provided at the other end of the tube, and the tube is filled with a gas. The rod is made to oscillate longitudinally at frequency \$f\$ to produce sound waves inside the gas, and the location of the plunger is adjusted until a standing sound wave pattern is set up inside the tube. Once the standing wave is set up, the motion of the gas molecules causes the cork filings to collect in a pattern of ridges at the displacement nodes.

\#\# Question

If \$f = 4.46 \textbackslash{}times 10\textasciicircum{}3\$ Hz and the separation between ridges is \$9.20\$ cm, what is the speed of sound in the gas?

\#\# Answer choices

- A: A: 756 m/s

- B: B: 778 m/s

- C: C: 793 m/s

- D: D: 821 m/s

\#\# Figure caption (auxiliary; use the image as primary evidence)

The image is a schematic diagram featuring several labeled components:

1. **Text Labels**:
   - "R": Positioned on the left, associated with a pipe or tube.
   - "D": Located near the center, next to a vertical structure.
   - "P": Positioned on the right, associated with another pipe or tube.
   - "d": In the middle between "D" and "P", with double-headed arrows indicating a distance.

2. **Objects**:
   - Pipes or tubes: Horizontal blue tubes located on both the left and right sides of the image.
   - Vertical section: A textured red rectangle above a gray base on the left.
   - Two vertical separators: Positioned near "D" and "P"; they are filled areas separating sections.
   - Particles or lumps: Brown shapes distributed between "D" and "P".

3. **Relationships**:
   - The pipes appear to connect through a narrow channel that includes the distance "d".
   - The particles or lumps are located between the two separators labeled "D" and "P" within the enclosed space. 

This illustrates a flow scenario or section of a system like a fluid or gas pipeline, possibly referencing filtration or separation processes.

\#\# Dataset metadata

Resonance and Harmonics; Physical Model Grounding Reasoning, Multi-Formula Reasoning

\paragraph{Recorded answer/context.}
D: D: 821 m/s

\paragraph{Figure/image path.}
/root/proposal\_for\_physic/science-mango/phyx\_mini\_run/phyx\_data/test\_image/328.png

\paragraph{Formalization target.}
create a compiling Lean file with sorry bodies at `PhyXMiniProblems/problem\_phyx\_mini\_0328.lean`. The Lean declarations must preserve the physical quantities, dimensions or dimensional roles, figure labels, governing-law hypotheses, and final relation expressed by this problem.
Use Mathlib/Physlib names found through LeanExplore where available. If a physics API is missing, introduce faithful local abstractions rather than scalar placeholder aliases.

\begin{theorem}[Physics formalization target]
\label{thm:physics:phyx_mini_0328:target}
\lean{PhyXMiniProblems.ProblemPhyXMini0328.problem_phyx_mini_0328}
\uses{def:physics:phyx-mini-0328:phyxminiproblems-problemphyxmini0328-meterspersecondvalue, def:physics:phyx-mini-0328:phyxminiproblems-problemphyxmini0328-resonantgastubesetup, def:physics:phyx-mini-0328:phyxminiproblems-problemphyxmini0328-matchesresonanttubescenarioandfigure, def:physics:phyx-mini-0328:phyxminiproblems-problemphyxmini0328-matchesresonanttubenumericalreadouts, def:physics:phyx-mini-0328:phyxminiproblems-problemphyxmini0328-hasphysicalresonanttubeparameters, def:physics:phyx-mini-0328:phyxminiproblems-problemphyxmini0328-satisfiescorkridgestandingwavelaw, def:physics:phyx-mini-0328:phyxminiproblems-problemphyxmini0328-satisfiesacousticwavespeedlaw, def:physics:phyx-mini-0328:phyxminiproblems-problemphyxmini0328-recordeddatasetanswer, def:physics:phyx-mini-0328:phyxminiproblems-problemphyxmini0328-isuniqueroundeddisplayedanswer}
The assigned autoformalize agent should translate this physics problem into Lean declarations in the covered file, with theorem and lemma proof bodies written as `by sorry`.
\end{theorem}
\begin{proof}
This is an autoformalization task, not a proof task. Produce statements that can later be proved by the physics prover without weakening the physical model.
\end{proof}

% --- Archon declaration topology begin ---
\section{Declaration topology}
\begin{definition}[Acoustic Length]
  \label{def:physics:phyx-mini-0328:phyxminiproblems-problemphyxmini0328-acousticlength}
  \lean{PhyXMiniProblems.ProblemPhyXMini0328.AcousticLength}
  A nonnegative physical length, independent of the selected length unit
\end{definition}
\begin{proof}
  This is the declared model definition of Acoustic Length.
\end{proof}
\begin{definition}[Acoustic Frequency]
  \label{def:physics:phyx-mini-0328:phyxminiproblems-problemphyxmini0328-acousticfrequency}
  \lean{PhyXMiniProblems.ProblemPhyXMini0328.AcousticFrequency}
  A nonnegative physical frequency, carrying inverse-time dimension
\end{definition}
\begin{proof}
  This is the declared model definition of Acoustic Frequency.
\end{proof}
\begin{definition}[length Readout]
  \label{def:physics:phyx-mini-0328:phyxminiproblems-problemphyxmini0328-lengthreadout}
  \lean{PhyXMiniProblems.ProblemPhyXMini0328.lengthReadout}
  \uses{def:physics:phyx-mini-0328:phyxminiproblems-problemphyxmini0328-acousticlength}
  Read a physical length as a real scalar in the selected length unit
\end{definition}
\begin{proof}
  This is the declared model definition of length Readout.
\end{proof}
\begin{definition}[frequency Readout]
  \label{def:physics:phyx-mini-0328:phyxminiproblems-problemphyxmini0328-frequencyreadout}
  \lean{PhyXMiniProblems.ProblemPhyXMini0328.frequencyReadout}
  \uses{def:physics:phyx-mini-0328:phyxminiproblems-problemphyxmini0328-acousticfrequency}
  Read a physical frequency in inverse units of the selected time unit
\end{definition}
\begin{proof}
  This is the declared model definition of frequency Readout.
\end{proof}
\begin{definition}[speed Readout]
  \label{def:physics:phyx-mini-0328:phyxminiproblems-problemphyxmini0328-speedreadout}
  \lean{PhyXMiniProblems.ProblemPhyXMini0328.speedReadout}
  Read a physical speed in the selected length unit per selected time unit
\end{definition}
\begin{proof}
  This is the declared model definition of speed Readout.
\end{proof}
\begin{definition}[meters Value]
  \label{def:physics:phyx-mini-0328:phyxminiproblems-problemphyxmini0328-metersvalue}
  \lean{PhyXMiniProblems.ProblemPhyXMini0328.metersValue}
  \uses{def:physics:phyx-mini-0328:phyxminiproblems-problemphyxmini0328-acousticlength, def:physics:phyx-mini-0328:phyxminiproblems-problemphyxmini0328-lengthreadout}
  Metre readout used for the standing wavelength and tube geometry
\end{definition}
\begin{proof}
  This is the declared model definition of meters Value.
\end{proof}
\begin{definition}[centimeters Value]
  \label{def:physics:phyx-mini-0328:phyxminiproblems-problemphyxmini0328-centimetersvalue}
  \lean{PhyXMiniProblems.ProblemPhyXMini0328.centimetersValue}
  \uses{def:physics:phyx-mini-0328:phyxminiproblems-problemphyxmini0328-acousticlength, def:physics:phyx-mini-0328:phyxminiproblems-problemphyxmini0328-lengthreadout}
  Centimetre readout used for the stated ridge separation
\end{definition}
\begin{proof}
  This is the declared model definition of centimeters Value.
\end{proof}
\begin{definition}[hertz Value]
  \label{def:physics:phyx-mini-0328:phyxminiproblems-problemphyxmini0328-hertzvalue}
  \lean{PhyXMiniProblems.ProblemPhyXMini0328.hertzValue}
  \uses{def:physics:phyx-mini-0328:phyxminiproblems-problemphyxmini0328-acousticfrequency, def:physics:phyx-mini-0328:phyxminiproblems-problemphyxmini0328-frequencyreadout}
  Hertz readout of the rod's cyclic driving frequency
\end{definition}
\begin{proof}
  This is the declared model definition of hertz Value.
\end{proof}
\begin{definition}[meters Per Second Value]
  \label{def:physics:phyx-mini-0328:phyxminiproblems-problemphyxmini0328-meterspersecondvalue}
  \lean{PhyXMiniProblems.ProblemPhyXMini0328.metersPerSecondValue}
  \uses{def:physics:phyx-mini-0328:phyxminiproblems-problemphyxmini0328-speedreadout}
  Metres-per-second readout used by the requested answer
\end{definition}
\begin{proof}
  This is the declared model definition of meters Per Second Value.
\end{proof}
\begin{definition}[Figure Label]
  \label{def:physics:phyx-mini-0328:phyxminiproblems-problemphyxmini0328-figurelabel}
  \lean{PhyXMiniProblems.ProblemPhyXMini0328.FigureLabel}
  Literal labels printed in the primary figure
\end{definition}
\begin{proof}
  This is the declared model definition of Figure Label.
\end{proof}
\begin{definition}[Figure Component]
  \label{def:physics:phyx-mini-0328:phyxminiproblems-problemphyxmini0328-figurecomponent}
  \lean{PhyXMiniProblems.ProblemPhyXMini0328.FigureComponent}
  Physical or geometric object denoted by each printed figure label
\end{definition}
\begin{proof}
  This is the declared model definition of Figure Component.
\end{proof}
\begin{definition}[Rod Support]
  \label{def:physics:phyx-mini-0328:phyxminiproblems-problemphyxmini0328-rodsupport}
  \lean{PhyXMiniProblems.ProblemPhyXMini0328.RodSupport}
  How the oscillating rod is supported
\end{definition}
\begin{proof}
  This is the declared model definition of Rod Support.
\end{proof}
\begin{definition}[Rod Oscillation]
  \label{def:physics:phyx-mini-0328:phyxminiproblems-problemphyxmini0328-rodoscillation}
  \lean{PhyXMiniProblems.ProblemPhyXMini0328.RodOscillation}
  Direction of the rod motion that excites the gas
\end{definition}
\begin{proof}
  This is the declared model definition of Rod Oscillation.
\end{proof}
\begin{definition}[Tube Material]
  \label{def:physics:phyx-mini-0328:phyxminiproblems-problemphyxmini0328-tubematerial}
  \lean{PhyXMiniProblems.ProblemPhyXMini0328.TubeMaterial}
  Material of the containing tube
\end{definition}
\begin{proof}
  This is the declared model definition of Tube Material.
\end{proof}
\begin{definition}[Tracer Material]
  \label{def:physics:phyx-mini-0328:phyxminiproblems-problemphyxmini0328-tracermaterial}
  \lean{PhyXMiniProblems.ProblemPhyXMini0328.TracerMaterial}
  Material distributed inside the tube to reveal displacement nodes
\end{definition}
\begin{proof}
  This is the declared model definition of Tracer Material.
\end{proof}
\begin{definition}[Acoustic Medium]
  \label{def:physics:phyx-mini-0328:phyxminiproblems-problemphyxmini0328-acousticmedium}
  \lean{PhyXMiniProblems.ProblemPhyXMini0328.AcousticMedium}
  Kind of acoustic medium occupying the resonant tube
\end{definition}
\begin{proof}
  This is the declared model definition of Acoustic Medium.
\end{proof}
\begin{definition}[Gas Column End]
  \label{def:physics:phyx-mini-0328:phyxminiproblems-problemphyxmini0328-gascolumnend}
  \lean{PhyXMiniProblems.ProblemPhyXMini0328.GasColumnEnd}
  The two axial ends of the resonant gas column
\end{definition}
\begin{proof}
  This is the declared model definition of Gas Column End.
\end{proof}
\begin{definition}[Boundary Apparatus]
  \label{def:physics:phyx-mini-0328:phyxminiproblems-problemphyxmini0328-boundaryapparatus}
  \lean{PhyXMiniProblems.ProblemPhyXMini0328.BoundaryApparatus}
  Boundary apparatus at either axial end of the gas column
\end{definition}
\begin{proof}
  This is the declared model definition of Boundary Apparatus.
\end{proof}
\begin{definition}[Cork Ridge Label]
  \label{def:physics:phyx-mini-0328:phyxminiproblems-problemphyxmini0328-corkridgelabel}
  \lean{PhyXMiniProblems.ProblemPhyXMini0328.CorkRidgeLabel}
  The four cork ridges visible from left to right in the primary figure
\end{definition}
\begin{proof}
  This is the declared model definition of Cork Ridge Label.
\end{proof}
\begin{definition}[Adjacent Cork Ridges]
  \label{def:physics:phyx-mini-0328:phyxminiproblems-problemphyxmini0328-adjacentcorkridges}
  \lean{PhyXMiniProblems.ProblemPhyXMini0328.AdjacentCorkRidges}
  \uses{def:physics:phyx-mini-0328:phyxminiproblems-problemphyxmini0328-corkridgelabel}
  Consecutive cork ridges in the left-to-right order shown in the figure
\end{definition}
\begin{proof}
  This is the declared model definition of Adjacent Cork Ridges.
\end{proof}
\begin{definition}[Resonant Gas Tube Setup]
  \label{def:physics:phyx-mini-0328:phyxminiproblems-problemphyxmini0328-resonantgastubesetup}
  \lean{PhyXMiniProblems.ProblemPhyXMini0328.ResonantGasTubeSetup}
  \uses{def:physics:phyx-mini-0328:phyxminiproblems-problemphyxmini0328-acousticlength, def:physics:phyx-mini-0328:phyxminiproblems-problemphyxmini0328-acousticfrequency, def:physics:phyx-mini-0328:phyxminiproblems-problemphyxmini0328-figurelabel, def:physics:phyx-mini-0328:phyxminiproblems-problemphyxmini0328-figurecomponent, def:physics:phyx-mini-0328:phyxminiproblems-problemphyxmini0328-rodsupport, def:physics:phyx-mini-0328:phyxminiproblems-problemphyxmini0328-rodoscillation, def:physics:phyx-mini-0328:phyxminiproblems-problemphyxmini0328-tubematerial, def:physics:phyx-mini-0328:phyxminiproblems-problemphyxmini0328-tracermaterial, def:physics:phyx-mini-0328:phyxminiproblems-problemphyxmini0328-acousticmedium, def:physics:phyx-mini-0328:phyxminiproblems-problemphyxmini0328-gascolumnend, def:physics:phyx-mini-0328:phyxminiproblems-problemphyxmini0328-boundaryapparatus, def:physics:phyx-mini-0328:phyxminiproblems-problemphyxmini0328-corkridgelabel}
  The physical resonant-tube setup. The sound speed and wavelength are independent fields: neither is assigned the answer value here
\end{definition}
\begin{proof}
  This is the declared model definition of Resonant Gas Tube Setup.
\end{proof}
\begin{definition}[Cork Ridges Are Strictly Ordered]
  \label{def:physics:phyx-mini-0328:phyxminiproblems-problemphyxmini0328-corkridgesarestrictlyordered}
  \lean{PhyXMiniProblems.ProblemPhyXMini0328.CorkRidgesAreStrictlyOrdered}
  \uses{def:physics:phyx-mini-0328:phyxminiproblems-problemphyxmini0328-metersvalue, def:physics:phyx-mini-0328:phyxminiproblems-problemphyxmini0328-resonantgastubesetup}
  The depicted ridges occur in strict left-to-right axial order
\end{definition}
\begin{proof}
  This is the declared model definition of Cork Ridges Are Strictly Ordered.
\end{proof}
\begin{definition}[Matches Resonant Tube Scenario And Figure]
  \label{def:physics:phyx-mini-0328:phyxminiproblems-problemphyxmini0328-matchesresonanttubescenarioandfigure}
  \lean{PhyXMiniProblems.ProblemPhyXMini0328.MatchesResonantTubeScenarioAndFigure}
  \uses{def:physics:phyx-mini-0328:phyxminiproblems-problemphyxmini0328-lengthreadout, def:physics:phyx-mini-0328:phyxminiproblems-problemphyxmini0328-metersvalue, def:physics:phyx-mini-0328:phyxminiproblems-problemphyxmini0328-adjacentcorkridges, def:physics:phyx-mini-0328:phyxminiproblems-problemphyxmini0328-resonantgastubesetup, def:physics:phyx-mini-0328:phyxminiproblems-problemphyxmini0328-corkridgesarestrictlyordered}
  Qualitative apparatus information and geometry read from the problem and the primary figure. The d marker spans the two central adjacent ridges. Its numerical 9.20 cm readout is recorded separately
\end{definition}
\begin{proof}
  This is the declared model definition of Matches Resonant Tube Scenario And Figure.
\end{proof}
\begin{definition}[Matches Resonant Tube Numerical Readouts]
  \label{def:physics:phyx-mini-0328:phyxminiproblems-problemphyxmini0328-matchesresonanttubenumericalreadouts}
  \lean{PhyXMiniProblems.ProblemPhyXMini0328.MatchesResonantTubeNumericalReadouts}
  \uses{def:physics:phyx-mini-0328:phyxminiproblems-problemphyxmini0328-centimetersvalue, def:physics:phyx-mini-0328:phyxminiproblems-problemphyxmini0328-hertzvalue, def:physics:phyx-mini-0328:phyxminiproblems-problemphyxmini0328-resonantgastubesetup}
  Numerical measurements stated in the problem
\end{definition}
\begin{proof}
  This is the declared model definition of Matches Resonant Tube Numerical Readouts.
\end{proof}
\begin{definition}[Has Physical Resonant Tube Parameters]
  \label{def:physics:phyx-mini-0328:phyxminiproblems-problemphyxmini0328-hasphysicalresonanttubeparameters}
  \lean{PhyXMiniProblems.ProblemPhyXMini0328.HasPhysicalResonantTubeParameters}
  \uses{def:physics:phyx-mini-0328:phyxminiproblems-problemphyxmini0328-metersvalue, def:physics:phyx-mini-0328:phyxminiproblems-problemphyxmini0328-hertzvalue, def:physics:phyx-mini-0328:phyxminiproblems-problemphyxmini0328-meterspersecondvalue, def:physics:phyx-mini-0328:phyxminiproblems-problemphyxmini0328-resonantgastubesetup}
  Positivity and containment conditions for the physical apparatus
\end{definition}
\begin{proof}
  This is the declared model definition of Has Physical Resonant Tube Parameters.
\end{proof}
\begin{definition}[Satisfies Cork Ridge Standing Wave Law]
  \label{def:physics:phyx-mini-0328:phyxminiproblems-problemphyxmini0328-satisfiescorkridgestandingwavelaw}
  \lean{PhyXMiniProblems.ProblemPhyXMini0328.SatisfiesCorkRidgeStandingWaveLaw}
  \uses{def:physics:phyx-mini-0328:phyxminiproblems-problemphyxmini0328-lengthreadout, def:physics:phyx-mini-0328:phyxminiproblems-problemphyxmini0328-corkridgelabel, def:physics:phyx-mini-0328:phyxminiproblems-problemphyxmini0328-adjacentcorkridges, def:physics:phyx-mini-0328:phyxminiproblems-problemphyxmini0328-resonantgastubesetup}
  Cork-ridge observation law for this standing sound wave. The ridges mark displacement nodes, and consecutive displacement nodes are separated by one half-wavelength. This is a general standing-wave law and does not assign a numerical wavelength or speed
\end{definition}
\begin{proof}
  This is the declared model definition of Satisfies Cork Ridge Standing Wave Law.
\end{proof}
\begin{definition}[Satisfies Acoustic Wave Speed Law]
  \label{def:physics:phyx-mini-0328:phyxminiproblems-problemphyxmini0328-satisfiesacousticwavespeedlaw}
  \lean{PhyXMiniProblems.ProblemPhyXMini0328.SatisfiesAcousticWaveSpeedLaw}
  \uses{def:physics:phyx-mini-0328:phyxminiproblems-problemphyxmini0328-lengthreadout, def:physics:phyx-mini-0328:phyxminiproblems-problemphyxmini0328-frequencyreadout, def:physics:phyx-mini-0328:phyxminiproblems-problemphyxmini0328-speedreadout, def:physics:phyx-mini-0328:phyxminiproblems-problemphyxmini0328-resonantgastubesetup}
  Nondispersive acoustic propagation law v = λ f for the driven gas wave, required in every compatible selection of length and time units
\end{definition}
\begin{proof}
  This is the declared model definition of Satisfies Acoustic Wave Speed Law.
\end{proof}
\begin{lemma}[marked Ridge Separation is half Wavelength]
  \label{lem:physics:phyx-mini-0328:phyxminiproblems-problemphyxmini0328-markedridgeseparation-is-halfwavelength}
  \lean{PhyXMiniProblems.ProblemPhyXMini0328.markedRidgeSeparation_is_halfWavelength}
  \uses{def:physics:phyx-mini-0328:phyxminiproblems-problemphyxmini0328-metersvalue, def:physics:phyx-mini-0328:phyxminiproblems-problemphyxmini0328-resonantgastubesetup, def:physics:phyx-mini-0328:phyxminiproblems-problemphyxmini0328-matchesresonanttubescenarioandfigure, def:physics:phyx-mini-0328:phyxminiproblems-problemphyxmini0328-satisfiescorkridgestandingwavelaw}
  The marked adjacent-ridge separation is one half-wavelength
\end{lemma}
\begin{proof}
  The conclusion follows from the governing relations and figure readouts named in the statement of marked Ridge Separation is half Wavelength.
\end{proof}
\begin{lemma}[standing Wavelength In Gas meters eq]
  \label{lem:physics:phyx-mini-0328:phyxminiproblems-problemphyxmini0328-standingwavelengthingas-meters-eq}
  \lean{PhyXMiniProblems.ProblemPhyXMini0328.standingWavelengthInGas_meters_eq}
  \uses{def:physics:phyx-mini-0328:phyxminiproblems-problemphyxmini0328-metersvalue, def:physics:phyx-mini-0328:phyxminiproblems-problemphyxmini0328-resonantgastubesetup, def:physics:phyx-mini-0328:phyxminiproblems-problemphyxmini0328-matchesresonanttubescenarioandfigure, def:physics:phyx-mini-0328:phyxminiproblems-problemphyxmini0328-matchesresonanttubenumericalreadouts, def:physics:phyx-mini-0328:phyxminiproblems-problemphyxmini0328-hasphysicalresonanttubeparameters, def:physics:phyx-mini-0328:phyxminiproblems-problemphyxmini0328-satisfiescorkridgestandingwavelaw}
  A 9.20 cm node spacing determines a wavelength of 0.184 m
\end{lemma}
\begin{proof}
  The conclusion follows from the governing relations and figure readouts named in the statement of standing Wavelength In Gas meters eq.
\end{proof}
\begin{definition}[Answer Choice]
  \label{def:physics:phyx-mini-0328:phyxminiproblems-problemphyxmini0328-answerchoice}
  \lean{PhyXMiniProblems.ProblemPhyXMini0328.AnswerChoice}
  Labels of the four answer choices printed with the problem
\end{definition}
\begin{proof}
  This is the declared model definition of Answer Choice.
\end{proof}
\begin{definition}[meters Per Second]
  \label{def:physics:phyx-mini-0328:phyxminiproblems-problemphyxmini0328-answerchoice-meterspersecond}
  \lean{PhyXMiniProblems.ProblemPhyXMini0328.AnswerChoice.metersPerSecond}
  \uses{def:physics:phyx-mini-0328:phyxminiproblems-problemphyxmini0328-answerchoice}
  Metres-per-second value displayed beside each answer label
\end{definition}
\begin{proof}
  This is the declared model definition of meters Per Second.
\end{proof}
\begin{definition}[recorded Dataset Answer]
  \label{def:physics:phyx-mini-0328:phyxminiproblems-problemphyxmini0328-recordeddatasetanswer}
  \lean{PhyXMiniProblems.ProblemPhyXMini0328.recordedDatasetAnswer}
  \uses{def:physics:phyx-mini-0328:phyxminiproblems-problemphyxmini0328-answerchoice}
  The answer label recorded by the supplied dataset metadata
\end{definition}
\begin{proof}
  This is the declared model definition of recorded Dataset Answer.
\end{proof}
\begin{definition}[Rounds To Displayed Speed]
  \label{def:physics:phyx-mini-0328:phyxminiproblems-problemphyxmini0328-roundstodisplayedspeed}
  \lean{PhyXMiniProblems.ProblemPhyXMini0328.RoundsToDisplayedSpeed}
  \uses{def:physics:phyx-mini-0328:phyxminiproblems-problemphyxmini0328-meterspersecondvalue, def:physics:phyx-mini-0328:phyxminiproblems-problemphyxmini0328-resonantgastubesetup, def:physics:phyx-mini-0328:phyxminiproblems-problemphyxmini0328-answerchoice}
  Agreement after rounding a physical speed to the nearest whole m/s
\end{definition}
\begin{proof}
  This is the declared model definition of Rounds To Displayed Speed.
\end{proof}
\begin{definition}[Is Unique Rounded Displayed Answer]
  \label{def:physics:phyx-mini-0328:phyxminiproblems-problemphyxmini0328-isuniqueroundeddisplayedanswer}
  \lean{PhyXMiniProblems.ProblemPhyXMini0328.IsUniqueRoundedDisplayedAnswer}
  \uses{def:physics:phyx-mini-0328:phyxminiproblems-problemphyxmini0328-resonantgastubesetup, def:physics:phyx-mini-0328:phyxminiproblems-problemphyxmini0328-answerchoice, def:physics:phyx-mini-0328:phyxminiproblems-problemphyxmini0328-roundstodisplayedspeed}
  A choice is the unique displayed value to which the physical speed rounds
\end{definition}
\begin{proof}
  This is the declared model definition of Is Unique Rounded Displayed Answer.
\end{proof}
% --- Archon declaration topology end ---
% --- Archon physics formalization source end ---
